% GENERATED FILE. Edit canonical lessons, then run npm run generate:books.
% canonical-source-hash: fnv1a64:2912293f816257f9
% canonical-lessons: PT-C286-delito, PT-C286-vandalo, PT-C286-vandalismo, PT-C286-vitima, PT-C286-homicidio

\chapter{Delitos e vítimas}
\label{ch:pt-delitos-e-vitimas}
\hlchaptermodality{286}

\begin{chapteropening}
\textbf{By the end of this chapter:} \emph{I can name offences, vandals and vandalism, victims and a homicide.}

Complete the last lesson of chapter 286: I can name offences, vandals and vandalism, victims and a homicide.
\end{chapteropening}

\section[o delito]{o delito — the offence, the misdeed}

\label{lesson:PT-C286-delito}



\begin{quote}
Before the new one: say the Portuguese for the TV viewer, then the Portuguese for the reader.
\end{quote}

\subsection*{You'll want to know: o delito}
\textbf{o delito} — \textquotedblleft{}the offence, the misdeed\textquotedblright{}.

\begin{quote}
Conduzir sem carta é um delito.
\end{quote}

\begin{quote}
\emph{Driving without a licence is an offence.}
\end{quote}

\emph{Em flagrante delito} — caught in the act.

\begin{cousinweb}
\emph{delito} was borrowed from Latin \emph{dēlictum} \textquotedblleft{}a fault, a misdeed\textquotedblright{} (EtymDB), from \emph{dēlinquere} \textquotedblleft{}to fall short, to do wrong\textquotedblright{}.
\end{cousinweb}

\begin{grammarlens}[title={What you've built}]
One more word for telling, explaining and describing at B1.
\end{grammarlens}

\subsection*{Your turn}
\begin{itemize}
\raggedright
  \item \emph{Say it:} \emph{o delito}
  \item \emph{Say it:} \emph{o delito}, once more
  \item \emph{Say it:} a sentence of your own with \emph{o delito}
  \item \emph{Recall:} say the Portuguese for the TV viewer, then the Portuguese for the reader, then say \emph{o delito} again
\end{itemize}

\begin{tcolorbox}[breakable,colback=teal!4,colframe=teal!35!black,title={Before you move on}]
What does o delito mean? (\textquotedblleft{}The offence, the misdeed\textquotedblright{}.) Say it once more.
\end{tcolorbox}
\section[o vândalo]{o vândalo — the vandal}

\label{lesson:PT-C286-vandalo}



\begin{quote}
Before the new one: say the Portuguese for the reader, then the Portuguese for the offence.
\end{quote}

\subsection*{You'll want to know: o vândalo}
\textbf{o vândalo} — \textquotedblleft{}the vandal\textquotedblright{}.

\begin{quote}
Um vândalo pintou a estátua do parque.
\end{quote}

\begin{quote}
\emph{A vandal painted the statue in the park.}
\end{quote}

Named after an ancient Germanic people.

\begin{cousinweb}
\emph{vândalo} is a learned borrowing from Latin \emph{Vandalus} (EtymDB), the name of an ancient Germanic people.
\end{cousinweb}

\begin{grammarlens}[title={What you've built}]
One more word for telling, explaining and describing at B1.
\end{grammarlens}

\subsection*{Your turn}
\begin{itemize}
\raggedright
  \item \emph{Say it:} \emph{o vândalo}
  \item \emph{Say it:} \emph{o vândalo}, once more
  \item \emph{Say it:} a sentence of your own with \emph{o vândalo}
  \item \emph{Recall:} say the Portuguese for the reader, then the Portuguese for the offence, then say \emph{o vândalo} again
\end{itemize}

\begin{tcolorbox}[breakable,colback=teal!4,colframe=teal!35!black,title={Before you move on}]
What does o vândalo mean? (\textquotedblleft{}The vandal\textquotedblright{}.) Say it once more.
\end{tcolorbox}
\section[o vandalismo]{o vandalismo — vandalism}

\label{lesson:PT-C286-vandalismo}



\begin{quote}
Before the new one: say the Portuguese for the offence, then the Portuguese for the vandal.
\end{quote}

\subsection*{You'll want to know: o vandalismo}
\textbf{o vandalismo} — \textquotedblleft{}vandalism\textquotedblright{}.

\begin{quote}
O vandalismo na estação custou muito dinheiro.
\end{quote}

\begin{quote}
\emph{The vandalism at the station cost a lot of money.}
\end{quote}

\emph{Atos de vandalismo} — acts of vandalism.

\begin{cousinweb}
\emph{vandalismo} is built on \emph{vândalo} with \emph{-ismo}, an ending that comes from Greek \emph{-ismós} through Latin (EtymDB has \emph{vândalo}).
\end{cousinweb}

\begin{grammarlens}[title={What you've built}]
One more word for telling, explaining and describing at B1.
\end{grammarlens}

\subsection*{Your turn}
\begin{itemize}
\raggedright
  \item \emph{Say it:} \emph{o vandalismo}
  \item \emph{Say it:} \emph{o vandalismo}, once more
  \item \emph{Say it:} a sentence of your own with \emph{o vandalismo}
  \item \emph{Recall:} say the Portuguese for the offence, then the Portuguese for the vandal, then say \emph{o vandalismo} again
\end{itemize}

\begin{tcolorbox}[breakable,colback=teal!4,colframe=teal!35!black,title={Before you move on}]
What does o vandalismo mean? (\textquotedblleft{}Vandalism\textquotedblright{}.) Say it once more.
\end{tcolorbox}
\section[a vítima]{a vítima — the victim}

\label{lesson:PT-C286-vitima}



\begin{quote}
Before the new one: say the Portuguese for the vandal, then the Portuguese for vandalism.
\end{quote}

\subsection*{You'll want to know: a vítima}
\textbf{a vítima} — \textquotedblleft{}the victim\textquotedblright{}.

\begin{quote}
A vítima foi levada para o hospital.
\end{quote}

\begin{quote}
\emph{The victim was taken to hospital.}
\end{quote}

Feminine in form for a man or a woman: \emph{ele foi a vítima}.

\begin{cousinweb}
\emph{vítima} is a learned borrowing from Latin \emph{victima} \textquotedblleft{}an animal offered in sacrifice\textquotedblright{} (Lewis \& Short).
\end{cousinweb}

\begin{grammarlens}[title={What you've built}]
One more word for telling, explaining and describing at B1.
\end{grammarlens}

\subsection*{Your turn}
\begin{itemize}
\raggedright
  \item \emph{Say it:} \emph{a vítima}
  \item \emph{Say it:} \emph{a vítima}, once more
  \item \emph{Say it:} a sentence of your own with \emph{a vítima}
  \item \emph{Recall:} say the Portuguese for the vandal, then the Portuguese for vandalism, then say \emph{a vítima} again
\end{itemize}

\begin{tcolorbox}[breakable,colback=teal!4,colframe=teal!35!black,title={Before you move on}]
What does a vítima mean? (\textquotedblleft{}The victim\textquotedblright{}.) Say it once more.
\end{tcolorbox}
\section[o homicídio]{o homicídio — the homicide, the killing}

\label{lesson:PT-C286-homicidio}



\begin{quote}
Before the new one: say the Portuguese for vandalism, then the Portuguese for the victim.
\end{quote}

\subsection*{You'll want to know: o homicídio}
\textbf{o homicídio} — \textquotedblleft{}the homicide, the killing\textquotedblright{}.

\begin{quote}
A polícia acha que foi um homicídio.
\end{quote}

\begin{quote}
\emph{The police think it was a homicide.}
\end{quote}

A news and court word; in speech, \emph{o assassinato}.

\begin{cousinweb}
\emph{homicídio} is a learned borrowing from Latin \emph{homicīdium} \textquotedblleft{}a killing of a person\textquotedblright{}, from \emph{homō} \textquotedblleft{}a human being\textquotedblright{} and \emph{caedere} \textquotedblleft{}to cut down, to kill\textquotedblright{} (Lewis \& Short).
\end{cousinweb}

\begin{grammarlens}[title={What you've built}]
One more word for telling, explaining and describing at B1.
\end{grammarlens}

\subsection*{Your turn}
\begin{itemize}
\raggedright
  \item \emph{Say it:} \emph{o homicídio}
  \item \emph{Say it:} \emph{o homicídio}, once more
  \item \emph{Say it:} a sentence of your own with \emph{o homicídio}
  \item \emph{Recall:} say the Portuguese for vandalism, then the Portuguese for the victim, then say \emph{o homicídio} again
\end{itemize}

\begin{tcolorbox}[breakable,colback=teal!4,colframe=teal!35!black,title={Before you move on}]
What does o homicídio mean? (\textquotedblleft{}The homicide, the killing\textquotedblright{}.) Say it once more.
\end{tcolorbox}
